% =============================================================================
\section{Calificar, depurar y chuleta}
% =============================================================================

\begin{frame}{Criterio de éxito de cada misión}
  \centering\scriptsize
  \begin{tabular}{clp{8.2cm}}
    \toprule
    \# & Misión & El nodo imprime ``Misión cumplida'' si\ldots \\
    \midrule
    1 & Estrella & pasa a \textless 0.5 m de cada vértice, regresa al centro y aterriza \\
    2 & Pilar & radio $3 \pm 0.3$ m, yaw \textless 10°, $\geq$ 90\% de las vueltas y llega a 4 m \\
    3 & Entrega & secuencia completa sin tocar QGC y termina a \textless 1.5 m de la H de la farmacia \\
    4 & Geocerca & el dron se detiene antes de pasar 1 m la cerca (prueben 2 y 6 m/s, varias direcciones) \\
    5 & Cueva & cruza, el lidar nunca a \textless 0.6 m de la roca, aterriza a \textless 1 m de la plataforma \\
    6 & Caza & a la vista $\geq$ 80\% de 60 s y a \textless 1.5 m en promedio \\
    7 & Estación & cada carga aterriza a \textless 0.3 m del centro del ArUco y termina la patrulla \\
    \bottomrule
  \end{tabular}
\end{frame}

\begin{frame}{La rúbrica}
  \centering\small
  \begin{tabular}{lrp{7cm}}
    \toprule
    Criterio & Peso & Qué se revisa \\
    \midrule
    Funciona & 40\% & ``Misión cumplida'' en 3 corridas seguidas \\
    Manejo de fallas & 15\% & Aborta y regresa bien cuando se le provoca una falla \\
    Código legible & 20\% & Máquina de estados clara, comentarios como en \cmd{px4\_examples} \\
    Parámetros & 10\% & Se ajusta con \cmd{--ros-args} sin recompilar \\
    README + video & 15\% & Otra persona lo corre solo con el README \\
    \bottomrule
  \end{tabular}
  \vspace{0.6em}
  \begin{idea}
    \small Cada quien presenta su misión en 5 minutos frente al equipo. Explicar la máquina de estados en voz
    alta es donde se nota si de verdad la entendieron.
  \end{idea}
\end{frame}

\begin{frame}{Provocar fallas a propósito}
  \small El 15\% de ``manejo de fallas'' se prueba así:
  \begin{itemize}
    \item \textbf{Matar el nodo} a media misión (Ctrl+C): en offboard PX4 debe entrar en failsafe;
      con comandos (misión 3) se queda en Hold.
    \item \textbf{Tomar el control desde QGC} (cambiar a Position o Hold): el nodo debe soltar y avisar
      ``Fuera de OFFBOARD''.
    \item \textbf{Geocerca}: escaparse en distintas direcciones y velocidades con \cmd{vuelo\_prueba}.
    \item \textbf{Cueva}: subir \cmd{max\_speed} o bajar \cmd{stop\_distance} hasta que falle; ¿cómo falla?
    \item \textbf{Estación}: subir \cmd{low\_percent} a 0.9 para que vaya a cargar más seguido.
  \end{itemize}
\end{frame}

\begin{frame}{Problemas comunes (1/2)}
  \centering\scriptsize
  \begin{tabular}{p{4.4cm}p{8.4cm}}
    \toprule
    Síntoma & Causa y arreglo \\
    \midrule
    No arma; \cmd{No connection to the GCS} & QGC no está abierto: \cmd{vespo qgc} (o usar \cmd{vespo up --sim}) \\
    \cmd{PX4 rechazó el comando 176} & El modo pedido no se puede ahora (p. ej. armar desde Land): pasar por Hold \\
    El dron se va al suelo o hacia arriba & Signo de $z$: en NED arriba es negativo \\
    La trayectoria se adelanta al dron & Tiempo con \cmd{now()} y la simulación lenta: usar \cmd{px4\_.clock()} \\
    ``Fuera de OFFBOARD'' sin tocar nada & El nodo dejó de publicar \textgreater0.5 s (un \cmd{return} antes de \cmd{offboard()}, un bloqueo) \\
    \bottomrule
  \end{tabular}
\end{frame}

\begin{frame}{Problemas comunes (2/2)}
  \centering\scriptsize
  \begin{tabular}{p{4.4cm}p{8.4cm}}
    \toprule
    Síntoma & Causa y arreglo \\
    \midrule
    \cmd{/x500/scan} o \cmd{/x500/camera/image} vacíos & El puente no está: usar el \emph{launch}; revisar el modelo (\cmd{--custom}) y el nombre del mundo \\
    El detector no ve nada & Falta \cmd{camera\_info} (el detector la espera), marcador muy chico o sin margen blanco \\
    La RAM de Gazebo crece sin parar & Cámara pesada o plugin de video de PX4: usar \cmd{x500\_cam\_abajo} y nuestros mundos \\
    La laptop se apaga sola & Sin cargador: conectar \\
    Un mundo modificado no cambia & Faltó \cmd{build\_worlds.py} y \cmd{vespo build} \\
    \bottomrule
  \end{tabular}
\end{frame}

\begin{frame}[fragile]{Chuleta}
  \begin{lstlisting}[style=bash,caption={Lo de todos los días}]
vespo build                                            # compilar
vespo up --sim --world <mundo> [--custom <modelo>]      # simular
ros2 run x500_retos <nodo> --ros-args -p <param>:=<valor>
ros2 launch x500_retos {cueva,caza,estacion}.launch.py
ros2 topic echo /fmu/out/vehicle_status_v1 --field nav_state
ros2 run rqt_image_view rqt_image_view /x500/aruco/debug
python3 src/x500_retos/tools/build_worlds.py            # regenerar mundos
gz topic -e -t /world/<mundo>/stats -n 1 | grep real_time   # ¿va a 1.0x?
  \end{lstlisting}
\end{frame}

\begin{frame}{Para cerrar}
  \begin{columns}[c]
    \column{0.55\textwidth}\small
    \begin{enumerate}
      \item \textbf{Nivel 1}: muévanse bien (geometría, derivadas, comandos).
      \item \textbf{Nivel 2}: que el dron decida (predecir, reaccionar a un sensor).
      \item \textbf{Nivel 3}: que vea (percepción separada del control).
    \end{enumerate}
    \vspace{0.4em}
    No se salta de nivel sin cerrar al menos dos misiones del anterior.
    \column{0.42\textwidth}\centering
    \includegraphics[width=0.75\linewidth]{vespo_logo}\\[0.5em]
    {\Large ¿Preguntas?}
  \end{columns}
\end{frame}
